\scanpage{70}
d'autre part $(a)$ s'écrit, modulo \uline{shift} de degré et tensorisation par \deleted{\unsure{la $\Omega^d$}} un $\underline\Omega^d_{X'/Y'}$
(NB $\underline\Omega^d_{X'/Y'}\simeq\underline\Omega^d_{X/Y}\otimes_{\mathcal O_X}\mathcal O_{X'}$)
\[
g'^*(\mathbb{R}i^!)\simeq(\mathbb{R}j^!)g^*
\]
et résulte donc simplement du \uline{th. de changement de base lisse}.

Cela définit l'isom de foncteurs qu'on voulait définir. Il reste à vérifier la compatibilité avec les hom de dualité, ce qui n'offre pas de difficultés, et est laissé au lecteur \unsure{\ldots}

\uline{Remarques}\quad a) Pour se sentir à l'aise, il aurait fallu vérifier que les isom ainsi définis entre les foncteurs $\mathbb{R}^!f$, par des immersions différentes dans des $X'$ lisses $/Y$, satisfont la condition naturelle de transitivité.
Nous \unsure{l'admettrons|l'admettons} --- la vérification ne \unsure{saurait} être particulièrement triviale. Cette compatibilité et de nombreuses autres seront \unsure{énoncées} plus loin dans le ``\uline{formulaire}''.
\mspage{69}
